\NormalFont\ShortTitle{2 Tesaloniceni}

{\MT A doua scrisoare a lui Pavel către Tesaloniceni


\par }\ChapOne{1}{\PP \VerseOne{1}Pavel, Silvan și Timotei, către adunarea Tesalonicenilor, în Dumnezeu, Tatăl nostru, și în Domnul Isus Hristos:

\VS{2}Har vouă și pace de la Dumnezeu, Tatăl nostru, și de la Domnul Isus Hristos.


\par }{\PP \VS{3}Suntem datori să mulțumim totdeauna lui Dumnezeu pentru voi, fraților, așa cum se cuvine, pentru că credința voastră crește nespus de mult și dragostea fiecăruia dintre voi unul față de altul este din ce în ce mai mare,

\VS{4}astfel încât noi înșine ne lăudăm cu voi în adunările lui Dumnezeu pentru perseverența și credința voastră în toate prigoanele și necazurile pe care le îndurați.

\VS{5}Acesta este un semn evident al dreptei judecăți a lui Dumnezeu, pentru ca voi să fiți socotiți vrednici de Împărăția lui Dumnezeu, pentru care și voi suferiți.

\VS{6}Căci este un lucru drept la Dumnezeu să răsplătească suferința celor care vă chinuiesc pe voi,

\VS{7}și să vă ușureze pe voi, cei care sunteți chinuiți împreună cu noi, când Domnul Isus se va arăta din cer cu îngerii săi puternici în flăcări de foc,

\VS{8}pedepsindu-i pe cei care nu-l cunosc pe Dumnezeu și pe cei care nu ascultă de Buna Vestire a Domnului nostru Isus,

\VS{9}care va plăti pedeapsa: distrugerea veșnică de la fața Domnului și de la gloria puterii sale,

\VS{10}când va veni în ziua aceea ca să fie glorificat în sfinții săi și să fie admirat printre toți cei care au crezut, pentru că mărturia noastră către voi a fost crezută.


\par }{\PP \VS{11}Pentru aceasta ne rugăm și noi totdeauna pentru voi, ca Dumnezeul nostru să vă socotească vrednici de chemarea voastră și să împlinească cu putere orice dorință de bunătate și orice lucrare de credință,

\VS{12}pentru ca Numele Domnului nostru Isus să fie proslăvit în voi și voi în El, după harul Dumnezeului nostru și al Domnului Isus Hristos.



\par }\Chap{2}{\PP \VerseOne{1}Acum, fraților, cu privire la venirea Domnului nostru Isus Hristos și la adunarea noastră la El, vă rugăm

\VS{2}să nu vă cutremurați repede în mintea voastră și să nu vă tulburați, nici prin duh, nici prin cuvânt, nici prin scrisoare, ca și cum ar veni de la noi, spunând că ziua lui Hristos a venit deja.

\VS{3}Nimeni să nu vă înșele în vreun fel. Căci aceasta nu va fi decât dacă va veni mai întâi răzvrătirea și dacă se va arăta omul păcatului, fiul pierzării.

\VS{4}El se opune și se înalță împotriva a tot ceea ce se numește Dumnezeu sau la care se închină, astfel încât se așează ca Dumnezeu în templul lui Dumnezeu, dându-se drept Dumnezeu.

\VS{5}Nu vă amintiți că, atunci când eram încă cu voi, v-am spus aceste lucruri?

\VS{6}Acum știți ce-l oprește, ca să se arate la vremea lui.

\VS{7}Căci misterul nelegiuirii deja lucrează. Numai că este unul care îl reține acum, până când va fi scos din cale.

\VS{8}Atunci se va arăta cel fără de lege, pe care Domnul îl va ucide cu suflarea gurii Sale și îl va nimici prin manifestarea venirii Sale;

\VS{9}cel a cărui venire este conform lucrării lui Satana, cu toată puterea, cu semne și minuni mincinoase,

\VS{10}și cu toată înșelăciunea răutății pentru cei care se pierd, pentru că nu au primit dragostea adevărului, ca să fie mântuiți.

\VS{11}Din această cauză, Dumnezeu le trimite o amăgire puternică, ca să creadă o minciună,

\VS{12}pentru ca să fie judecați toți cei care nu au crezut adevărul, ci au avut plăcere în nedreptate.


\par }{\PP \VS{13}Dar noi suntem datori să mulțumim totdeauna lui Dumnezeu pentru voi, frați iubiți de Domnul, pentru că Dumnezeu v-a ales de la început pentru mântuirea prin sfințirea Duhului și prin credința în adevăr,

\VS{14}la care v-a chemat prin Buna Vestire, în vederea obținerii slavei Domnului nostru Isus Hristos.

\VS{15}Așadar, fraților, rămâneți deci tari și păstrați tradițiile pe care ați fost învățați de noi, fie prin cuvânt, fie prin scrisoare.


\par }{\PP \VS{16}Iar Domnul nostru Iisus Hristos însuși și Dumnezeu Tatăl nostru, care ne-a iubit și ne-a dat prin har mângâierea veșnică și buna nădejde,

\VS{17}să mângâie inimile voastre și să vă întărească în orice faptă bună și în orice cuvânt.



\par }\Chap{3}{\PP \VerseOne{1}În sfârșit, fraților, rugați-vă pentru noi, ca să se răspândească repede Cuvântul Domnului și să fie proslăvit, ca și la voi,

\VS{2}și să fim izbăviți de oamenii nechibzuiți și răi, căci nu toți au credință.

\VS{3}Dar credincios este Domnul, care vă va întări și vă va păzi de cel rău.

\VS{4}Noi avem încredere în Domnul în ceea ce vă privește, că voi faceți și veți face ceea ce vă poruncim.

\VS{5}Domnul să vă îndrepte inimile spre dragostea lui Dumnezeu și spre perseverența lui Hristos.


\par }{\PP \VS{6}Și acum, fraților, vă poruncim, în numele Domnului nostru Isus Hristos, să vă îndepărtați de orice frate care umblă în răzvrătire și nu după tradiția pe care au primit-o de la noi.

\VS{7}Căci voi știți cum trebuie să ne imitați. Căci noi nu ne-am purtat cu răzvrătire între voi,

\VS{8}și nici nu am mâncat pâine din mâna cuiva fără să o plătim, ci, în muncă și osteneală, am lucrat zi și noapte, ca să nu împovărăm pe niciunul dintre voi.

\VS{9}Aceasta nu pentru că nu avem dreptul, ci pentru a ne face noi înșine un exemplu pentru voi, ca să ne imitați.

\VS{10}Căci și când eram cu voi, v-am poruncit acest lucru: “Dacă cineva nu vrea să muncească, să nu-l lăsați să mănânce”.

\VS{11}Căci am auzit că unii dintre voi umblă cu răzvrătire, care nu muncesc deloc, ci sunt niște oameni ocupați.

\VS{12}Or, celor care sunt așa, le poruncim și îi îndemnăm în Domnul Isus Hristos să lucreze cu liniște și să-și mănânce propria pâine.


\par }{\PP \VS{13}Dar voi, fraților, nu vă obosiți în a face ce este drept.

\VS{14}Dacă cineva nu se supune cuvântului nostru din această scrisoare, observați-l pe acel om și nu vă apropiați de el, ca să fie rușinat.

\VS{15}Nu-l socotiți ca pe un dușman, ci admonestați-l ca pe un frate.


\par }{\PP \VS{16}Și Domnul păcii să vă dea pace în toate zilele și în toate felurile. Domnul să fie cu voi toți.


\par }{\PP \VS{17}Eu, Pavel, scriu cu mâna mea această salutare, care este semnul din fiecare scrisoare. Iată cum scriu eu.

\VS{18}Harul Domnului nostru Isus Hristos să fie cu voi toți. Amin.


\par }